\documentclass[11pt]{article}
\usepackage[margin=1in]{geometry}
\usepackage{amsmath,amssymb,amsthm}
\usepackage{graphicx}
\usepackage{booktabs}
\usepackage{enumitem}
\usepackage[hidelinks]{hyperref}

\newtheorem{theorem}{Theorem}
\newtheorem{lemma}{Lemma}
\newcommand{\cstar}{c_0^*}
\newcommand{\ellstar}{(\ell^1)^*}
\newcommand{\ainf}{a_\infty}
\newcommand{\schar}{\Sigma}
\newcommand{\F}{F}
\newcommand{\se}{\operatorname{se}}

\title{The Arithmetic-Mean-at-Infinity Cancellation Envelope\\
of Every Principal Operator Ideal Is Finite Rank}
\author{Candidate full solution packet for arXiv:0707.3271, Question 7}
\date{9 August 2026}

\begin{document}
\maketitle

\begin{center}
\fbox{\textbf{Status: full solution, likely valid; pending specialist review}}
\end{center}

\section{Source problem}

Kaftal and Weiss define, for an operator ideal $I$,
\[
  \widehat I^{\,\infty}
  :=\bigcap\{J: J_{\ainf}\supset I_{\ainf}\},
\]
and ask the following in Question 7 of their survey \cite{KWsurvey}, on PDF
page 24.

\begin{quote}
If $I=(\xi)$, is $\widehat I^{\,\infty}$ principal? If yes, is there an
expression for its generator similar to the one in Definition 5.15 and with
similar properties?
\end{quote}

\begin{figure}[h]
  \centering
  \includegraphics[width=\textwidth]{figures/open_problem_crop.png}
  \caption{Question 7 and its defining context in \cite[p.~24]{KWsurvey}.}
\end{figure}

\section{Notation and source facts}

Let $\cstar$ be the cone of nonnegative nonincreasing null sequences and let
$\ellstar$ be its summable part.  For $m\in\mathbb N$, the ampliation $D_m$
repeats every entry $m$ times, so
\[
 (D_m\eta)_n=\eta_{\lceil n/m\rceil}.
\]
For $\eta\in\ellstar$, write
\[
  (\eta_{\ainf})_n=\frac1n\sum_{j=n+1}^{\infty}\eta_j.
\]
If $(\eta)$ is the principal ideal generated by $\eta$, Calkin's
characteristic-set criterion, in the notation of the source, is
\begin{equation}\label{eq:principal-membership}
 \rho\in\schar((\eta))
 \quad\Longleftrightarrow\quad
 \rho\le C D_m\eta\quad\text{for some }C>0,\ m\in\mathbb N.
\end{equation}

We use the following identity of Kaftal and Weiss
\cite[Lemma~4.7]{KWtraces}.  For an infinite-rank $\xi\in\cstar$,
\begin{equation}\label{eq:KW-lemma}
 (\xi)_{\ainf}=
 \begin{cases}
   (\xi_{\ainf}),&\xi\in\ell^1,\\
   \se(\omega),&\xi\notin\ell^1,
 \end{cases}
\end{equation}
where $\omega=\langle 1/n\rangle$.  The finite-rank case will be handled
directly.  Finally, $\F$ denotes the ideal of finite-rank operators, the least
nonzero ideal of $B(H)$.

\section{Proof intuition}

The key observation is that domination by an arithmetic mean at infinity is a
tail-mass condition, not a pointwise lower bound.  A fixed amount of tail mass
may be spread over an arbitrarily long constant block, forcing that block's
height to be as small as desired.

Given any infinite-rank test sequence $\rho$, we exploit this freedom
diagonally.  In the summable case we replace consecutive blocks of the source
generator $\xi$ by their averages.  This preserves each block's mass and
increases every intermediate tail, but the blocks can be chosen long enough
that their heights defeat every fixed ampliation $D_m$ in
\eqref{eq:principal-membership}.  In the nonsummable case, long low plateaux
with mass at least one per block give a nonsummable generator, so
\eqref{eq:KW-lemma} makes its arithmetic-mean-at-infinity ideal equal to the
same universal ideal $\se(\omega)$; the same diagonal argument excludes
$\rho$.  Thus every infinite-rank sequence is absent from at least one ideal
in the defining intersection.

\section{The classification}

\begin{theorem}\label{thm:main}
For every nonzero principal ideal $I$ of $B(H)$,
\[
  \boxed{\widehat I^{\,\infty}=\F.}
\]
Consequently $\widehat I^{\,\infty}$ is principal, independently of the
generator of $I$, and any nonzero finite-rank operator is a generator.
\end{theorem}

\begin{proof}
First, $I_{\ainf}\ne\{0\}$ for nonzero principal $I$.  Hence every $J$ in
the defining intersection for $\widehat I^{\,\infty}$ is nonzero.  Since
$\F$ is contained in every nonzero operator ideal,
\begin{equation}\label{eq:lower}
  \F\subset\widehat I^{\,\infty}.
\end{equation}
We prove the reverse inclusion by showing that each infinite-rank
$\rho\in\cstar$ is omitted by some qualifying principal ideal $J$.  Because
$\rho$ is nonincreasing and has infinite rank, $\rho_n>0$ for every $n$.

\smallskip
\noindent\textbf{Case 1: $I=\F$.}
We first note that $\F_{\ainf}=\F$.  One inclusion is immediate because the
arithmetic mean at infinity of a finite-support sequence again has finite
support.  Conversely, if $\alpha$ is supported on $\{1,\dots,r\}$, take
$\zeta=\langle\underbrace{1,\dots,1}_{r+1},0,\dots\rangle$.  Then
$(\zeta_{\ainf})_n=(r+1-n)/n>0$ for $1\le n\le r$, so
$\alpha\le C\zeta_{\ainf}$ for a suitable $C$.  Thus
$\alpha\in\schar(\F_{\ainf})$.  It follows that $J=\F$ itself belongs to the
defining family, and therefore $\widehat{\F}^{\,\infty}=\F$.

\smallskip
\noindent\textbf{Case 2: $I=(\xi)$, where $\xi\in\ellstar$ has infinite
rank.}
Set $N_0=0$.  Inductively, once $N_{k-1}$ is chosen, put
$q_k=N_{k-1}+1$ and choose $N_k>N_{k-1}$ so large that
\begin{equation}\label{eq:hk-summable}
 h_k:=\frac{\sum_{j=q_k}^{N_k}\xi_j}{N_k-N_{k-1}}
       <\frac{\rho_{kq_k}}{k}.
\end{equation}
This is possible because the numerator is bounded by $\|\xi\|_1$ while the
denominator may be arbitrarily large.  Define
\[
 \eta_j=h_k\qquad(q_k\le j\le N_k).
\]
The sequence $\eta$ is nonincreasing: the average on a later block is no
larger than its first $\xi$-entry, which is no larger than the last
$\xi$-entry of the preceding block, which in turn is no larger than the
preceding block average.  Moreover, block masses are preserved, so
$\eta\in\ellstar$ and $\sum_j\eta_j=\sum_j\xi_j$.

For $N_{k-1}\le n\le N_k$, the suffix
$\xi_{n+1},\dots,\xi_{N_k}$ has average no larger than the average of the
whole $k$-th block, because $\xi$ is nonincreasing.  Therefore
\[
 \sum_{j=n+1}^{N_k}\eta_j
   =(N_k-n)h_k
   \ge \sum_{j=n+1}^{N_k}\xi_j.
\]
All later block masses agree exactly.  Hence, for every $n$,
\begin{equation}\label{eq:tail-dominance}
 \sum_{j>n}\eta_j\ge\sum_{j>n}\xi_j,
 \qquad\text{and thus}\qquad
 \eta_{\ainf}\ge\xi_{\ainf}.
\end{equation}
Let $J=(\eta)$.  From \eqref{eq:KW-lemma} and
\eqref{eq:tail-dominance},
\[
 J_{\ainf}=(\eta_{\ainf})\supset(\xi_{\ainf})=I_{\ainf},
\]
so $J$ qualifies for the defining intersection.

We claim that $\rho\notin\schar(J)$.  Otherwise, by
\eqref{eq:principal-membership}, $\rho\le C D_m\eta$ for some $C>0$ and
$m\in\mathbb N$.  Choose an integer $k\ge m$ with $k>C$ and evaluate at
$n=mq_k$.  Since $(D_m\eta)_{mq_k}=\eta_{q_k}=h_k$, monotonicity of $\rho$
and \eqref{eq:hk-summable} give
\[
 \rho_{mq_k}\ge\rho_{kq_k}>kh_k
       >C(D_m\eta)_{mq_k},
\]
contradicting the proposed bound.

\smallskip
\noindent\textbf{Case 3: $I=(\xi)$, where $\xi\notin\ell^1$.}
By \eqref{eq:KW-lemma}, $I_{\ainf}=\se(\omega)$.  We build a nonsummable
$\eta\in\cstar$ whose principal ideal omits $\rho$.  Set $N_0=0$.  Having
chosen $N_{k-1}$, put $q_k=N_{k-1}+1$ and choose $h_k>0$ so that
\begin{equation}\label{eq:hk-nonsummable}
 h_k<\min\left\{\frac1k,\frac{h_{k-1}}2,
                      \frac{\rho_{kq_k}}k\right\}
\end{equation}
for $k>1$, omitting the middle term for $k=1$.  Let
$L_k=\lceil1/h_k\rceil$, set $N_k=N_{k-1}+L_k$, and put
$\eta_j=h_k$ for $q_k\le j\le N_k$.  Then $\eta$ is nonincreasing and tends
to zero, but
\[
 \sum_j\eta_j=\sum_k L_kh_k\ge\sum_k1=\infty.
\]
Thus $J=(\eta)$ is not contained in the trace class, and
\eqref{eq:KW-lemma} gives
\[
 J_{\ainf}=\se(\omega)=I_{\ainf}.
\]
It qualifies.  If $\rho\le C D_m\eta$, choose $k\ge m$ with $k>C$ and
again evaluate at $n=mq_k$.  Equation \eqref{eq:hk-nonsummable} yields
\[
 \rho_{mq_k}\ge\rho_{kq_k}>kh_k
       >C(D_m\eta)_{mq_k},
\]
the same contradiction.

We have shown that every infinite-rank $\rho\in\cstar$ is missing from at
least one ideal in the defining intersection.  Thus
$\widehat I^{\,\infty}\subset\F$, and \eqref{eq:lower} proves equality.
\end{proof}

\section{Verification notes}

The attempt was re-read adversarially against the exact source conventions.
The current verdict is \textbf{likely valid} (confidence $93/100$).

\begin{center}
\begin{tabular}{p{0.29\textwidth}p{0.62\textwidth}}
\toprule
Check & Finding\\
\midrule
All generator cases & Finite-rank, summable infinite-rank, and nonsummable
cases are disjoint and exhaustive.\\
Tail domination & A suffix of a decreasing finite block has average at most
the full-block average; future block masses are preserved exactly.\\
Ampliations & For fixed $m,C$, a stage $k\ge m$, $k>C$ contradicts the bound
at $mq_k$.  This explicitly handles the dilation freedom in principal-ideal
membership.\\
Nonsummability & Each plateau has mass $L_kh_k\ge1$, while $h_k\to0$.\\
External input & Equation \eqref{eq:KW-lemma} is precisely
\cite[Lemma~4.7]{KWtraces}; the principal-membership convention agrees with
\cite[\S2.1]{KWsurvey}.\\
Computation & None used or needed.\\
\bottomrule
\end{tabular}
\end{center}

The most important specialist checks are the implication from
\eqref{eq:tail-dominance} to $I_{\ainf}\subset J_{\ainf}$, the application of
Lemma 4.7 in Case 3, and the direct finite-rank identity
$\F_{\ainf}=\F$.

\section{Novelty search and limitations}

On 9 August 2026, the lightweight registry, solution, attempt, and proof-gap
indexes for the run were searched for arXiv:0707.3271, the exact Question 7
wording, and the core arithmetic-mean-at-infinity cancellation terms.  The
local parsed-source corpus was searched for notation and phrase variants.
Bounded web searches used the exact question, variants of
``$\widehat I^{\infty}$,'' ``arithmetic mean at infinity cancellation,'' and
the authors' names.  They returned the 2007 Kaftal--Weiss multipaper series
and their later work on majorization, but no separate paper recording an
answer to Question 7.  This supports, but does not guarantee, novelty.

The theorem answers Question 7 exactly.  It does not address the survey's
other nine questions or cancellation envelopes built from other arithmetic
mean operations.  The proof uses the known identity \eqref{eq:KW-lemma}; it
does not reprove that substantial background result.

\enlargethispage{5\baselineskip}
\small
\begin{thebibliography}{9}

\bibitem{KWsurvey}
V. Kaftal and G. Weiss,
\emph{A survey on the interplay between arithmetic mean ideals, traces,
lattices of operator ideals, and an infinite Schur--Horn majorization
theorem}, arXiv:0707.3271 (2007).

\bibitem{KWtraces}
V. Kaftal and G. Weiss,
\emph{Traces on operator ideals and arithmetic means},
arXiv:0707.3169 (2007), especially Lemma 4.7.

\end{thebibliography}
\normalsize

\end{document}
